\chapter{Discussion}
\label{ch:discussion}

\section{Capabilities of AMR-Based Detection}
\label{sec:disc-capabilities}

% AMR absorbs surface variation: synonym predicates, different wordings, embedded
%   instructions in structured fields all unify under the same graph structure.
% Zero false positives under Rule1+24: predicate-level agreement is a precise signal.
% Injection-safe sub-calls: single-concept synonym/antonym queries cannot be exploited
%   because the attacker controls only one concept node — insufficient for injection.
% Authorization annotation cleanly separates delegated from injected actions within scope.

\section{Limitations}
\label{sec:disc-limitations}

\subsection{LLM Parsing Non-Determinism}
% The AMR graph for a given text depends on the LLM parser; output varies across calls.
% Injections phrased far from PropBank canonical frames (euphemisms, buried sentences,
%   unusual metaphors) may not produce a graph that triggers any rule.
% Mitigation directions: constrained decoding, deterministic rule-based parsers for
%   well-typed domains.

\subsection{Broad-Delegation Vulnerability}
% When the user explicitly delegates a task file, :auth t applies to all instructions
%   within it, including injected ones.
% This is not a logic error but a fundamental limit of user-level authorization.
% Requires intent verification beyond the current AMR policy model.

\subsection{Structural Ambiguity}
% Concept-agnostic graph matching (Rule5 / smatch) cannot distinguish financially
%   malicious from financially benign AMR trees.
% Meaningful graph-based detection in typed domains requires concept-sensitive matching,
%   which reintroduces the need for semantic knowledge already supplied by Rule1+24.

\section{Abstraction Mismatch Extension}
\label{sec:disc-abstraction}

% Problem: "don't give financial advice" vs "buy Tesla stock" — normal predicate
%   matching fails because "buy" does not appear in the policy.
% Proposed extension: query the LLM with the full policy phrase and the top AMR
%   concept node from the injection ("buy-01 :mode imperative").
%   The LLM evaluates whether the concept constitutes an instance of the policy.
%   The attacker controls only "buy-01"; a single-node query is injection-safe
%   by the same argument as the synonym sub-calls.
% This extends the detection class to abstract policies without sacrificing the
%   injection-safety property of the LLM sub-calls.
% Sketch the call interface; note this is future work.

\section{Future Directions}
\label{sec:disc-future}

\subsection{Information Flow Taint Tracking}
% CaMeL-style approach: label each value in the agent's context with a trust level
%   (user-authorised vs. tool-output-derived).
% Track how tainted values flow through tool calls and reasoning steps.
% Flag when a tainted value influences a privileged action (e.g. a financial transfer).
% This complements the AMR semantic check: taint catches injections that change
%   \emph{parameters} of authorized actions rather than invoking forbidden predicates.

\subsection{User-Defined Variable Binding}
% Allow the user to bind named variables at task start (e.g. MY\_IBAN = "DE89...").
% Tool outputs cannot overwrite these bindings.
% Any instruction that references a different value for the same role is flagged.
% This directly addresses the IBAN-swap attack class without requiring predicate matching.

\subsection{Expanded Attack Taxonomy and Benchmark}
% The current suite covers categories A–J; systematic coverage of indirect injection
%   vectors (image alt-text, PDF metadata, calendar descriptions) is future work.
% Differential testing: compare firewall behaviour on raw text vs. AMR-replaced
%   message history to isolate the contribution of each component.
